% ---------------------------------------------------------------------------
% Table A1 --- the three coincidences sketched in Figure 1 (labor-market axis).
% Split out of the single Appendix A table on the author's 2026-10-22 request
% (##AG line 54): the figure-side rows first, because Figure 1 precedes Table 4.
% Rows keep the author's order and wording; repairs are typographical only
% (eta_s -> eta_{s,i} in the corner row, "as a assuming" -> "as assuming",
% "labor qunatities" -> "labor quantities", stray "allocation" deleted).
% Columns, status labels and their reading are unchanged.
% ---------------------------------------------------------------------------
\begin{table}[htbp]
\centering
\footnotesize
\setlength{\tabcolsep}{5pt}
\renewcommand{\arraystretch}{1.25}
\begin{tabular}{@{}>{\raggedright\arraybackslash}p{3.3cm}
                  >{\raggedright\arraybackslash}p{5.5cm}
                  >{\raggedright\arraybackslash}p{2.1cm}
                  >{\raggedright\arraybackslash}p{2.8cm}@{}}
\toprule
Link & Mechanism & Status & Left unidentified \\
\midrule
Full reallocation $\Leftrightarrow$ standard CGE
& The converse case of no-reallocation with $\eta = 1$ exactly represents the allocation rule the standard CGE closure imposes.
& definitional
& ---  \\[3pt]
Standard CGE $\Leftrightarrow$ single-wage elastic supply
& With the real wage at its anchor, the elastic labor supply curve and the full-employment solution give the same equation at every supply elasticity and under every financing rule.
& derivable (see below)
& the single-wage supply elasticity \\[3pt]
No reallocation $\Leftrightarrow$ the rigid corner of the sectoral family
& An infinitely inelastic supply curve ($\eta_{s,i}=0$) with sectoral wages gives the same results as assuming sectorally frozen labor quantities ($\eta=0$), albeit with a different rationale.
& exact by construction
& ---  \\[3pt]
Elastic labor supply $\rightarrow$ employment adjustment
& Under demand-only shocks, as $\eta_{s,i}\to\infty$ the supply curves flatten, prices return to the baseline and the real wage to its anchor; the family converges to the employment-adjustment allocation except for the net external transfer, which adjusts there and is fixed at zero here.
& limit, attained at a different transfer
& the transfer $F$ \\
\bottomrule
\end{tabular}
\caption{Links among the labor-market closures that
Figure~\ref{fig:flexibility-map} sketches, the mechanism behind each, their
logical status, and the parameter each leaves unidentified. The last line is
the boundary case: a limit rather than an identity, and one that lands beside
the closure rather than on it.}
\label{tab:equiv-links}
\end{table}
